\documentclass[12pt,a4paper]{article}

% --- Metadata ---
\def\courseshortheader{HỌC MÁY ỨNG DỤNG --- PANDAS}
\def\coverbookline{TRÍ TUỆ NHÂN TẠO}
\def\covermaintitle{Pandas DataFrame UltraQuick Tutorial}
\def\coversubtitle{Dịch tiếng Việt từ Google ML Crash Course (Colab)}

% Shared preamble for nhom_5 (requires \courseshortheader before \input)
\usepackage[utf8]{inputenc}
\usepackage[vietnamese]{babel}
\usepackage{amsmath,amssymb,amsthm}
\usepackage{geometry}
\usepackage{enumitem}
\usepackage[most]{tcolorbox}
\usepackage{hyperref}
\usepackage{fancyhdr}
\usepackage{graphicx}
\usepackage{booktabs}
\usepackage{tikz}
\usetikzlibrary{calc,arrows.meta,decorations.pathreplacing,positioning}
\usepackage{eso-pic}
\usepackage{xcolor}

\geometry{a4paper, margin=2cm, bottom=2.5cm}
\hypersetup{colorlinks=true, linkcolor=blue!70!black, urlcolor=blue}
\setlist[itemize]{topsep=4pt, itemsep=2pt}
\setlist[enumerate]{topsep=4pt, itemsep=2pt}

\AddToShipoutPictureFG{%
  \AtPageCenter{%
    \begin{tikzpicture}[remember picture, overlay]
      \node[rotate=45, scale=1, opacity=0.06]{%
        \fontsize{60}{72}\selectfont\bfseries\sffamily Đặng Tấn Quí%
      };
    \end{tikzpicture}%
  }%
}

\pagestyle{fancy}
\fancyhf{}
\fancyhead[L]{\small\textit{\courseshortheader}}
\fancyhead[R]{\small\textit{Đặng Tấn Quí}}
\fancyfoot[C]{\thepage}
\fancyfoot[L]{\footnotesize\textcopyright\ Đặng Tấn Quí}
\fancyfoot[R]{\footnotesize SĐT: 0377\,122\,210}
\renewcommand{\headrulewidth}{0.4pt}
\renewcommand{\footrulewidth}{0.4pt}

\fancypagestyle{plain}{%
  \fancyhf{}
  \fancyfoot[C]{\thepage}
  \fancyfoot[L]{\footnotesize\textcopyright\ Đặng Tấn Quí}
  \fancyfoot[R]{\footnotesize SĐT: 0377\,122\,210}
  \renewcommand{\headrulewidth}{0pt}
  \renewcommand{\footrulewidth}{0.4pt}
}

\newtcolorbox{dongco}[1][]{%
  enhanced, breakable,
  colback=teal!4!white, colframe=teal!60!black,
  fonttitle=\bfseries, title={Động cơ học -- Tại sao cần khái niệm này?}, #1%
}
\newtcolorbox{ynghia}[1][]{%
  enhanced, breakable,
  colback=purple!3!white, colframe=purple!50!black,
  fonttitle=\bfseries, title={Ý nghĩa \& Trực giác}, #1%
}
\newtcolorbox{ungdung}[1][]{%
  enhanced, breakable,
  colback=olive!5!white, colframe=olive!60!black,
  fonttitle=\bfseries, title={Ứng dụng thực tế}, #1%
}
\newtcolorbox{dinhnghia}[1][]{%
  enhanced, breakable,
  colback=blue!3!white, colframe=blue!60!black,
  fonttitle=\bfseries, title={Định nghĩa}, #1%
}
\newtcolorbox{dinhly}[1][]{%
  enhanced, breakable,
  colback=green!3!white, colframe=green!50!black,
  fonttitle=\bfseries, title={Định lý}, #1%
}
\newtcolorbox{tinhchat}[1][]{%
  enhanced, breakable,
  colback=yellow!5!white, colframe=orange!50!black,
  fonttitle=\bfseries, title={Tính chất}, #1%
}
\newtcolorbox{phuongphap}[1][]{%
  enhanced, breakable,
  colback=gray!5!white, colframe=gray!50!black,
  fonttitle=\bfseries, title={Phương pháp}, #1%
}
\newtcolorbox{luuy}[1][]{%
  enhanced, breakable,
  colback=red!3!white, colframe=red!50!black,
  fonttitle=\bfseries, title={Lưu ý}, #1%
}
\newtcolorbox{congthuc}[1][]{%
  enhanced, breakable,
  colback=cyan!3!white, colframe=cyan!50!black,
  fonttitle=\bfseries, title={Công thức}, #1%
}
\newtcolorbox{vidu}[1][]{%
  enhanced, breakable,
  colback=violet!3!white, colframe=violet!50!black,
  fonttitle=\bfseries, title={Ví dụ}, #1%
}
\newtcolorbox{phanvidu}[1][]{%
  enhanced, breakable,
  colback=red!2!white, colframe=red!40!black,
  fonttitle=\bfseries, title={Phản ví dụ}, #1%
}
\newtcolorbox{tongket}[1][]{%
  enhanced, breakable,
  colback=blue!4!white, colframe=blue!70!black,
  fonttitle=\bfseries, title={Tổng kết chương}, #1%
}


\usepackage{listings}
\usepackage{upquote}
\usepackage[T1]{fontenc}

\lstdefinestyle{mlcode}{
  basicstyle=\ttfamily\small,
  breaklines=true,
  breakatwhitespace=false,
  columns=fullflexible,
  keepspaces=true,
  showstringspaces=false,
  frame=single,
  framerule=0.4pt,
  rulecolor=\color{black!25},
  backgroundcolor=\color{black!2},
  xleftmargin=0.2cm,
  xrightmargin=0.2cm
}
\lstset{style=mlcode}

\newcommand{\mlccurl}[1]{\url{#1}}

\begin{document}

\thispagestyle{empty}
\begin{tikzpicture}[remember picture, overlay]
  \fill[blue!80!black] (current page.south west) rectangle (current page.north east);
  \fill[blue!60!cyan, opacity=0.8]
    (current page.south west) -- (current page.north west) --
    ([xshift=-3cm]current page.north east) -- (current page.south east) -- cycle;
  \foreach \r/\o in {6/0.05, 4.5/0.08, 3/0.12} {
    \fill[white, opacity=\o] ([xshift=2cm, yshift=-2cm]current page.north east) circle (\r cm);
  }
  \draw[white, line width=2pt] ([yshift=-5cm]current page.north west) ++(2cm,0) -- ++(17cm,0);
  \draw[white, line width=2pt] ([yshift=-16cm]current page.north west) ++(2cm,0) -- ++(17cm,0);
  \node[anchor=west, white, font=\fontsize{28}{34}\bfseries\selectfont]
    at ([xshift=3cm, yshift=-7.5cm]current page.north west) {\coverbookline};
  \node[anchor=west, white, font=\fontsize{20}{26}\selectfont]
    at ([xshift=3cm, yshift=-9cm]current page.north west) {\covermaintitle};
  \node[anchor=west, white, font=\fontsize{16}{22}\selectfont]
    at ([xshift=3cm, yshift=-10.2cm]current page.north west) {\coversubtitle};
  \node[anchor=west, white, font=\Large\bfseries]
    at ([xshift=3cm, yshift=-13cm]current page.north west) {Biên soạn:};
  \node[anchor=west, yellow!80!white, font=\fontsize{24}{30}\bfseries\selectfont]
    at ([xshift=3cm, yshift=-14.5cm]current page.north west) {ĐẶNG TẤN QUÍ};
  \node[anchor=west, white!90, font=\large]
    at ([xshift=3cm, yshift=-18cm]current page.north west) {SĐT: 0377\,122\,210};
  \node[anchor=south, white!80, font=\small]
    at ([yshift=1.5cm]current page.south) {Tài liệu có bản quyền --- Cấm sao chép khi chưa được sự đồng ý của tác giả};
  \node[anchor=south east, white, font=\Large\bfseries]
    at ([xshift=-2cm, yshift=3cm]current page.south east) {\the\year};
\end{tikzpicture}

\newpage
\tableofcontents
\newpage

%==============================================================================
\section*{Lời nói đầu}
\addcontentsline{toc}{section}{Lời nói đầu}
%==============================================================================

\begin{ynghia}[title={Nguồn}]
Tài liệu dịch từ notebook \textit{Pandas DataFrame UltraQuick Tutorial} của Google Machine Learning Crash Course:
\mlccurl{https://colab.research.google.com/github/google/eng-edu/blob/main/ml/cc/exercises/pandas_dataframe_ultraquick_tutorial.ipynb}
\end{ynghia}

\begin{luuy}[title={target}]
Đây không phải tutorial DataFrame đầy đủ. target là giới thiệu \textbf{cực nhanh} các phần của pandas DataFrame cần để làm các bài tập Colab trong ML Crash Course.
\end{luuy}

\newpage

%==============================================================================
\section{Colab là gì?}
%==============================================================================

\begin{dinhnghia}[title={Colaboratory (Colab)}]
Machine Learning Crash Course dùng \textbf{Colaboratory (Colab)} cho tất cả bài tập lập trình.
Colab là một hiện thực của Google dựa trên \textbf{Jupyter Notebook}:
\mlccurl{https://jupyter.org/}.

Để tìm hiểu cách dùng Colab, xem \textit{Welcome to Colaboratory}:
\mlccurl{https://research.google.com/colaboratory}.
\end{dinhnghia}

%==============================================================================
\section{Giới thiệu DataFrame}
%==============================================================================

\begin{dongco}[title={DataFrame là gì?}]
Notebook này giới thiệu \textbf{DataFrame}, cấu trúc dữ liệu trung tâm trong pandas.
Một DataFrame tương tự một bảng tính (spreadsheet) trong bộ nhớ:
\begin{itemize}
  \item DataFrame lưu dữ liệu trong các \textbf{ô}.
  \item DataFrame có các \textbf{cột có tên} (thường là vậy) và các \textbf{hàng đánh số}.
\end{itemize}
\end{dongco}

%==============================================================================
\section{Import NumPy và pandas}
%==============================================================================

\begin{phuongphap}[title={Import module}]
Chạy cell sau để import NumPy và pandas.
\end{phuongphap}

\begin{lstlisting}
import numpy as np
import pandas as pd
\end{lstlisting}

%==============================================================================
\section{Tạo một DataFrame}
%==============================================================================

\begin{phuongphap}[title={Tạo DataFrame từ mảng NumPy}]
Cell dưới đây tạo một DataFrame đơn giản gồm 10 ô theo cấu trúc:
\begin{itemize}
  \item 5 hàng
  \item 2 cột: \texttt{temperature} và \texttt{activity}
\end{itemize}

Ta khởi tạo \texttt{pd.DataFrame} với 2 tham số:
\begin{itemize}
  \item Tham số thứ nhất là dữ liệu để điền vào các ô (ở đây dùng \texttt{np.array} tạo mảng $5\times 2$).
  \item Tham số thứ hai là danh sách tên cột.
\end{itemize}

\textbf{Lưu ý:} Không đổi tên/ghi đè biến trong cell này vì các cell sau dùng lại.
\end{phuongphap}

\begin{lstlisting}
# Create and populate a 5x2 NumPy array.
my_data = np.array([[0, 3], [10, 7], [20, 9], [30, 14], [40, 15]])

# Create a Python list that holds the names of the two columns.
my_column_names = ['temperature', 'activity']

# Create a DataFrame.
my_dataframe = pd.DataFrame(data=my_data, columns=my_column_names)

# Print the entire DataFrame
print(my_dataframe)
\end{lstlisting}

%==============================================================================
\section{Thêm cột mới vào DataFrame}
%==============================================================================

\begin{phuongphap}[title={Tạo cột mới bằng phép gán}]
Bạn có thể thêm cột mới vào DataFrame bằng cách gán giá trị cho một tên cột mới.
Ví dụ sau tạo cột thứ ba \texttt{adjusted} trong \texttt{my\_dataframe}.
\end{phuongphap}

\begin{lstlisting}
# Create a new column named adjusted.
my_dataframe["adjusted"] = my_dataframe["activity"] + 2

# Print the entire DataFrame
print(my_dataframe)
\end{lstlisting}

%==============================================================================
\section{Chọn một phần (subset) của DataFrame}
%==============================================================================

\begin{phuongphap}[title={Lấy hàng/cột/đoạn con}]
pandas cung cấp nhiều cách để lấy ra các hàng, cột, lát cắt (slice) hoặc từng ô trong DataFrame.
\end{phuongphap}

\begin{lstlisting}
print("Rows #0, #1, and #2:")
print(my_dataframe.head(3), '\n')

print("Row #2:")
print(my_dataframe.iloc[[2]], '\n')

print("Rows #1, #2, and #3:")
print(my_dataframe[1:4], '\n')

print("Column 'temperature':")
print(my_dataframe['temperature'])
\end{lstlisting}

%==============================================================================
\section{Bài tập 1: Tạo một DataFrame}
%==============================================================================

\begin{phuongphap}[title={Yêu cầu bài tập}]
Thực hiện các bước sau:
\begin{enumerate}
  \item Tạo một pandas DataFrame kích thước $3 \times 4$ (3 hàng, 4 cột), với tên cột lần lượt là
  \texttt{Eleanor}, \texttt{Chidi}, \texttt{Tahani}, \texttt{Jason}. Điền mỗi ô bằng một số nguyên ngẫu nhiên trong đoạn \([0,100]\) (bao gồm cả 0 và 100).
  \item In ra:
  \begin{itemize}
    \item Toàn bộ DataFrame,
    \item Giá trị ở ô hàng \#1 của cột \texttt{Eleanor}.
  \end{itemize}
  \item Tạo cột thứ năm tên \texttt{Janet}, có giá trị bằng tổng theo từng hàng của \texttt{Tahani} và \texttt{Jason}.
\end{enumerate}

Để làm bài này, bạn sẽ cần kiến thức NumPy cơ bản trong tài liệu \textit{NumPy UltraQuick Tutorial}.
\end{phuongphap}

\begin{lstlisting}
# Write your code here.
\end{lstlisting}

\begin{lstlisting}
# Create a Python list that holds the names of the four columns.
my_column_names = ['Eleanor', 'Chidi', 'Tahani', 'Jason']

# Create a 3x4 numpy array, each cell populated with a random integer.
my_data = np.random.randint(low=0, high=101, size=(3, 4))

# Create a DataFrame.
df = pd.DataFrame(data=my_data, columns=my_column_names)

# Print the entire DataFrame
print(df)

# Print the value in row #1 of the Eleanor column.
print("\nSecond row of the Eleanor column: %d\n" % df['Eleanor'][1])

# Create a column named Janet whose contents are the sum
# of two other columns.
df['Janet'] = df['Tahani'] + df['Jason']

# Print the enhanced DataFrame
print(df)
\end{lstlisting}

%==============================================================================
\section{Sao chép DataFrame (tùy chọn)}
%==============================================================================

\begin{dongco}[title={Vì sao cần phân biệt tham chiếu và sao chép?}]
Pandas có hai cách khác nhau để ``nhân đôi'' một DataFrame:
\begin{itemize}
  \item \textbf{Tham chiếu (referencing):} nếu bạn gán một DataFrame cho một biến mới, mọi thay đổi ở một bên sẽ phản ánh sang bên kia.
  \item \textbf{Sao chép (copying):} nếu bạn gọi \texttt{pd.DataFrame.copy}, bạn tạo ra một bản sao độc lập thực sự; thay đổi ở bản gốc sẽ không ảnh hưởng bản sao và ngược lại.
\end{itemize}
Sự khác nhau có thể khó thấy ngay, nhưng rất quan trọng.
\end{dongco}

\begin{lstlisting}
# Create a reference by assigning my_dataframe to a new variable.
print("Experiment with a reference:")
reference_to_df = df

# Print the starting value of a particular cell.
print("  Starting value of df: %d" % df['Jason'][1])
print("  Starting value of reference_to_df: %d\n" % reference_to_df['Jason'][1])

# Modify a cell in df.
df.at[1, 'Jason'] = df['Jason'][1] + 5
print("  Updated df: %d" % df['Jason'][1])
print("  Updated reference_to_df: %d\n\n" % reference_to_df['Jason'][1])

# Create a true copy of my_dataframe
print("Experiment with a true copy:")
copy_of_my_dataframe = my_dataframe.copy()

# Print the starting value of a particular cell.
print("  Starting value of my_dataframe: %d" % my_dataframe['activity'][1])
print("  Starting value of copy_of_my_dataframe: %d\n" % copy_of_my_dataframe['activity'][1])

# Modify a cell in df.
my_dataframe.at[1, 'activity'] = my_dataframe['activity'][1] + 3
print("  Updated my_dataframe: %d" % my_dataframe['activity'][1])
print("  copy_of_my_dataframe does not get updated: %d" % copy_of_my_dataframe['activity'][1])
\end{lstlisting}

\end{document}

